% SSSUP-FONTE-STORICA: fonti-locali/originali/2014-2020.zip
% Anno 2018, esercizio 2. Trasposizione dal sorgente LaTeX storico.
% Verificare editorialmente la trasposizione prima di considerarla definitiva.
\renewcommand{\Anno}{2018}
\renewcommand{\NumeroProblema}{2}

\renewcommand{\Testo}{
Nello stato di Santannaland la patente a punti, segue le seguenti regole:
\begin{itemize}
\item al tempo $0$ tutte le patenti hanno $20$ punti;
\item ad ogni infrazione la patente viene decurtata di un punto;
\item se la patente arriva a $0$ punti, la patente viene sospesa;
\item dopo un tempo $L$ dall'ultima infrazione, i punti vengono riportati al valore massimo che la persona ha avuto nell'arco della sua vita;
\item se dall'ultimo istante in cui i punti sono stati riportati al loro massimo o dal tempo $0$ è trascorso un tempo $L$ senza infrazioni, allora il massimo viene innalzato di $2$ punti.
\end{itemize}

Supponiamo che ogni persona commetta un'infrazione ogni tempo $T$. Per cercare di massimizzare i loro punti, $2$ amici, in maniera fraudolenta, scelgono di comunicare chi tra loro ha compiuto l'infrazione e a cui verranno decurtarti quindi i punti, nel modo per loro più conveniente.

\begin{enumerate}
\item[(a)] Assumendo che i due amici adottino la strategia per loro più conveniente, quali sono i valori di $T$ che permettono di non sospendere alcuna patente?
\item[(b)] Nel caso in cui gli amici che si mettono d'accordo sono $n$, quali sono i valori di $T$ che permettono di non sospendere alcuna patente?
\end{enumerate}
}

\renewcommand{\Risposta}{
% SSSUP-BOZZA-AI: Risposta; autore=ChatGPT; data=2026-11-12; stato=DA-VERIFICARE
% SSSUP-PROVENIENZA: ChatGPT (OpenAI), sessione ProblemiSSSUP del 2026-11-12; derivato dalla traccia canonica; revisione umana pendente.
Assumiamo infrazioni periodiche di periodo $T$, con le prime in $(0,T]$,
e che un recupero maturato nello stesso istante di un'infrazione
venga applicato prima della decurtazione.

(a) Per due amici la condizione esatta \`e $T\ge L/10$.

(b) Per $n$ amici, indipendentemente dalle fasi delle infrazioni,
\[
 T\ge\frac{L}{20-\lceil19/n\rceil}\quad\hbox{\`e necessario},
 \qquad
 T\ge\frac{L}{19-\lceil18/n\rceil}\quad\hbox{\`e sufficiente}.
\]
Se le infrazioni complessive sono equidistanti, a intervalli $T/n$,
la condizione esatta \`e
\[
 \boxed{T\ge\frac{nL}{19n-18}.}
\]
Per $n$ generale le fasi possono cambiare la risposta: il Testo non
le precisa, perci\`o l'ultima formula non \`e una risposta incondizionata
all'enunciato originale.
}
\renewcommand{\Aiutino}{
% SSSUP-BOZZA-AI: Aiutino; autore=ChatGPT; data=2026-11-12; stato=DA-VERIFICARE
% SSSUP-PROVENIENZA: ChatGPT (OpenAI), sessione ProblemiSSSUP del 2026-11-12; derivato dalla traccia canonica; revisione umana pendente.
Senza alcun recupero, una patente pu\`o assorbire al massimo $19$
infrazioni: la ventesima la sospende.
Provare ad attribuire blocchi di $19$ infrazioni a un amico alla volta,
ciclicamente. Per dimostrare l'impossibilit\`a, considerare il primo
istante in cui qualcuno recupera punti: chi ha assorbito tutte le
infrazioni durante il suo intervallo di riposo?
Nel caso generale occorre distinguere infrazioni simultanee e sfalsate.
}
\renewcommand{\SoluzioneUfficiale}{}
\renewcommand{\SoluzioneRedazionale}{
% SSSUP-BOZZA-AI: SoluzioneRedazionale; autore=ChatGPT; data=2026-11-12; stato=DA-VERIFICARE
% SSSUP-PROVENIENZA: ChatGPT (OpenAI), sessione ProblemiSSSUP del 2026-11-12; derivato dalla traccia canonica; revisione umana pendente.
\textbf{Modello temporale e ambiguit\`a.}
Supponiamo $L,T>0$. Ciascun amico commette un'infrazione a intervalli
regolari $T$, con la prima in $(0,T]$; ogni infrazione deve essere
attribuita a uno degli amici. Le fasi possono essere diverse o coincidere.
Quando un recupero e una decurtazione cadono nello stesso istante,
applichiamo prima il recupero. Questa convenzione \`e necessaria per
includere le uguaglianze nelle condizioni sufficienti.
L'enunciato non precisa n\'e le fasi n\'e tale precedenza.

Una patente parte da $20$ punti e non pu\`o mai arrivare a zero:
senza recuperi pu\`o dunque assorbire al massimo $19$ infrazioni.
Le strategie che seguono non hanno bisogno dei bonus di due punti:
utilizzano soltanto il ripristino a un massimo che \`e almeno $20$.
I bonus non invalidano neppure le dimostrazioni di impossibilit\`a,
perch\'e anche il primo bonus richiede un intervallo $L$ senza infrazioni.

\textbf{Una condizione necessaria per $n$ amici.}
Poniamo
\[
 k=\left\lfloor\frac{19(n-1)}n\right\rfloor+1
   =20-\left\lceil\frac{19}n\right\rceil.
\]
Supponiamo per assurdo $L>kT$ e che si riesca a evitare tutte le sospensioni.
Prima o poi deve avvenire un aumento di punti, altrimenti le scorte
iniziali si esauriscono. Sia $t_*$ il primo istante di un aumento.
La patente che ne beneficia non ha ricevuto infrazioni nell'intervallo
aperto $(t_*-L,t_*)$, e $t_*\ge L$.
Un intervallo di lunghezza maggiore di $kT$ contiene almeno $k$
infrazioni commesse da ciascun amico, dunque almeno $nk$ complessive.
Queste devono essere state attribuite alle altre $n-1$ patenti.
Prima di $t_*$ nessuna ha recuperato punti e ciascuna pu\`o aver
assorbito al massimo $19$ infrazioni. Ma
\[
 nk>19(n-1),
\]
contraddizione. \`E quindi necessario $L\le kT$.

\textbf{Una strategia sufficiente per $n$ amici.}
Ordiniamo tutte le infrazioni, anche quelle simultanee, in una successione
di istanti $t_1\le t_2\le\cdots$. Per gli eventi simultanei scegliamo un
ordine che si ripeta a ogni periodo. Allora
\[
 t_{j+n}=t_j+T.
\]
Attribuiamo le prime $19$ infrazioni al primo amico, le successive $19$
al secondo, e cos\`i via; dopo l'ultimo amico ricominciamo dal primo.
Fra l'ultima infrazione di un blocco e la prima del blocco successivo
della stessa persona ci sono
\[
 D=19(n-1)+1=19n-18
\]
intervalli fra eventi consecutivi. Pertanto il tempo di riposo \`e almeno
\[
 t_{j+D}-t_j\ge
 \left\lfloor\frac Dn\right\rfloor T
 =\left(19-\left\lceil\frac{18}n\right\rceil\right)T.
\]
Se quest'ultima quantit\`a \`e almeno $L$, la patente ha recuperato
il proprio massimo prima del nuovo blocco. Ogni blocco toglie soltanto
$19$ punti, quindi nessuna patente viene sospesa.

\textbf{(a) Due amici.}
Per $n=2$ i due interi delle condizioni precedenti sono entrambi $10$.
La condizione necessaria e quella sufficiente coincidono:
\[
 \boxed{T\ge L/10.}
\]
In particolare, attribuire alternativamente $19$ infrazioni a ciascuno
garantisce a ogni patente un intervallo di riposo esattamente $10T$,
qualunque sia lo sfasamento delle due sequenze periodiche.
Con $T<L/10$, prima del primo recupero l'altro amico dovrebbe invece
assorbire almeno $20$ infrazioni.

\textbf{(b) Una risposta esatta con infrazioni equidistanti.}
Se le fasi sono uniformemente sfalsate, tutte le infrazioni complessive
avvengono a intervalli $\delta=T/n$.
La strategia a blocchi di $19$ concede un riposo esattamente
$D\delta$, quindi funziona se $L\le D\delta$.
Viceversa, se $L>D\delta$, nell'intervallo che precede il primo aumento
di punti si verificano almeno $D$ infrazioni, tutte a carico delle altre
$n-1$ patenti, ma $D=19(n-1)+1>19(n-1)$.
Dunque
\[
 L\le\frac{19n-18}{n}T
 \quad\Longleftrightarrow\quad
 T\ge\frac{nL}{19n-18}
\]
\`e condizione necessaria e sufficiente in questo modello.
Il caso $n=1$ restituisce correttamente $T\ge L$.

\textbf{Perch\'e non si possono omettere le fasi.}
Prendiamo $n=19$ e $L=18{,}5T$.
Se tutte le infrazioni sono simultanee agli istanti $T,2T,\ldots$,
attribuiamo le $19$ infrazioni di ogni istante a un solo amico,
ruotando fra i $19$ amici. Ogni patente perde $19$ punti, poi riposa
$19T>L$ prima di essere usata di nuovo: la strategia funziona.
Se invece le infrazioni sono equidistanti, il massimo intervallo
compatibile con la condizione appena dimostrata \`e
\[
 \frac{343}{19}T<18{,}5T=L,
\]
e nessuna strategia pu\`o funzionare.
Le due situazioni rispettano entrambe l'ipotesi di un'infrazione ogni $T$
per persona. Il testo, senza ulteriori precisazioni, non determina
quindi una sola condizione esatta per tutti gli $n$.
}
\renewcommand{\SoluzioniAlternative}{}
\renewcommand{\Osservazioni}{}
\renewcommand{\Commenti}{}
